\section{Property Names}\label{sec:NamesOfProperties}
\bdq{Properties}\index{property} are more of a concept than an actual element of the \Java programming language. They are characteristics of a class\index{class} accessed via accessors\index{method!accessor method} and mutators\index{method!mutator method} (or \bsq{getters}\index{method!getter method} and \bsq{setters}\index{method!setter method} in the case of JavaBeans\index{JavaBean}), though they do not necessarily have to exist as actual attributes\index{attribute} (fields\index{field|see {attribute}}) of that class. So properties\index{property} can actually be computed values. 

The JavaBeans\index{JavaBeans} specification \autocite{ORACLE_DOC_JAVABEANS:Chapter7} defines properties like this:
\begin{quote}
``Properties\index{property} are discrete, named attributes of a JavaBean\index{JavaBeans} that can affect its appearance or its behaviour. For example, a GUI button might have a property\index{propterty} named \bsq{Label} that represents the text displayed in the button.

[…]

Properties\index{property} can have arbitrary types, including both built-in Java types such as \bsq{\lstinline|int|} and class\index{class} or interface\index{interface} types such as \bsq{\lstinline|java.awt.Color|}\index{java.awt!Color}.

[…]

Properties\index{property} are always accessed via method\index{method} calls on their owning object.

[…] properties\index{property} need not just be simple data fields, they can actually be computed values. Updates may have various programmatic side effects.''
\end{quote}

Other definitions vary in details, but in general, they are in line with this one.

Based on that, the name of property\index{property} should be a noun, enhanced by an adjective, or an adjective (alone or in combination with \lstinline|is|, \lstinline|has| or alike) when it describes a state.

Usually, \textit{camelCase} is used for the property\index{•} names, with the first letter in lower case.

Samples are:
\begin{itemize}[nosep]
	\item{\verb#name#}
	\item{\verb#title#}
	\item{\verb#isReleased#}
	\item{\verb#hasChildren#}
\end{itemize}

%==============================================================================
% End of file!!
%==============================================================================
